\chapter{25}

\section{Bern (CH), Samstag, 25.
August}

\emph{Zürich-Bern-Brig-Betten Talstation} war auf das Papier gedruckt.
Der untenstehende Preis irritierte ihn nach wie vor. Die Schweizerischen
Bundesbahnen waren ganz schön teuer. Damit hatte er nicht gerechnet. Den
letzten Fünfziger musste er aufheben. Die \emph{Bettmeralp-Bahn} würde
er sich sparen müssen.

Vielleicht war sein Plan zu gewagt. Er durfte nicht alles auf den Hof
setzen. Wenn die bereits einen \emph{Zusenn} gefunden hatten? Oder, wenn
die Leute unsympathische, bornierte Bauern waren?

Mit einer resoluten Handbewegung wischte er die Gedanken weg.
Schwarzmalen brachte ihn nicht weiter. Nach längerem Anstehen bestieg er
den Zug. Marco. Der Gedanke an seinen Freund und Arbeitgeber knallte
mitten in dieser dichten Menschenmenge plötzlich und schon wieder in
seinen Kopf und liess ihn beinahe die Beine versagen. In Bern würde er
auch ihn anrufen. Dringend. Er würde es tun.

Der Waggon war brechend voll. Die Menschen waren vorwiegend mit
Rucksäcken und Reisekoffern unterwegs.

Es gelang Louis, einen Platz in einem Viererabteil zu bekommen. Rucksack
und Gitarre schob er auf die Ablage über seinem Kopf. Der ältere,
schwere Mann neben ihm berührte ihn mit jedem seiner Atemzüge. Louis
wich mehrmals nach rechts aus. Es wurde anstrengend. Dann liess er es.

Louis verband sein Handy mit den Kopfhörern und steckte sie in die
Ohren. Mit dem guten Gefühl an den vergangenen Abend suchte er nach
Alboráns \emph{«Pasos de cero»} und startete den Song. Es war Zeit für
ein bisschen Leichtigkeit. In Erinnerung an das musikalische
Miteinander. An Vanessa und Meo. Und an Meos Vater. Sein Saxophonspiel
hatte Louis umgehauen.

Sie hatten zusammen geklungen, als würden sie nichts anderes tun, als
seit Jahren gemeinsam musizieren. Und Meos Mutter hatte mit den Fingern
den Takt dazu geschnippt.

\sectionbreak

«Deine Stimme gefällt mir am besten, wenn du spanisch singst,» hatte
Giorgia einmal zu ihm gesagt. Bis dahin hatte er sich keine Gedanken
darum gemacht. Vielleicht hatte sie recht. Vielleicht klang er am
ehrlichsten, wenn er spanisch sang. Giorgias Aussage hatte damals
widersprüchliche Gefühle in ihm ausgelöst.

Bis zur Einschulung war er der Überzeugung gewesen, alle Väter würden
spanisch sprechen. Louis erinnerte sich, wie er Maman eines Tages
gefragt hatte, warum die Väter seiner Freunde Französisch redeten.

«Weil sie Schweizer sind,» hatte sie geantwortet, «und wir hier in der
französischen Schweiz leben. Papa liebt seine Sprache und will keine
andere lernen.»

Dann hatte sie ihm übers Haar gestrichen und ihn angeschaut, wie sie
immer schaute, wenn sie ihn beruhigen wollte.

«Ist doch prima für dich und deine Schwester. Ihr habt die Chance,
mehrere Sprachen zu lernen.»

Sein kindlicher Stolz, wenn er und Vater zusammen redeten und ausserhalb
der Familie sie niemand verstehen konnte, war etappenweise erstickt.
Louis’ Enttäuschung über sein unerwartetes Verschwinden hatte
schliesslich sein Vertrauen vollends schwinden lassen.

Wenigstens die spanische Sprache hatte ihm sein Vater hinterlassen,
erklärte er Giorgia an einem Abend missmutig und auch, dass er mit ihm
abgeschlossen habe, und sie ihn nicht mehr nach ihm fragen solle. Er
erinnerte sich an Giorgias verblüfften Gesichtsausdruck. Und dankbar
daran, dass sie tatsächlich nichts mehr sagte.

\sectionbreak

Louis schloss die Augen. Er versuchte das Sax-Solo von Meos Vater in
sich aufspielen zu lassen. Unerwartet hatte der es eingeflochten, als
wäre es ein selbstverständlicher Teil des Songs. Seine Spielweise hatte
etwas Sehnsüchtiges, Berührendes gehabt.

Der Rhythmus der Musik zusammen mit den regelmässigen Zugbewegungen
schaukelten Louis hin und her. Er döste vor sich hin.

\qrblock{Pablo Alborán \emph{«Pasos de cero»}}{media/QR-Codes/023_Pasos_de_cero_Pablo_Alboràn_Spotify.png}

Bis ihn der Schwere anstupste. Louis fuhr hoch. Dann schaute er in ein
freundliches Gesicht. Die junge Bahnangestellte hatte einen dunkeln
Teint, lange Rastazöpfe und einen Knipser in der Hand. Louis zog rasch
sein Ticket unter dem Handy hervor und streckte es ihr entgegen. «Ihr
Ticket, bitte» und, «vielen Dank,» hörte er sie sagen. Ihre Stimme hörte
sich melodiös an.

Den direkten Zugsanschluss nach Brig liess Louis sausen. Zugunsten des
geplanten Gesprächs mit Marco. Er hatte es sich und damit auch Marco
versprochen.

Seine Suche nach passenden Worten begann beim Betreten der Rolltreppe
nach oben. Ungeordnet noch. Er brauchte Tageslicht.

\sectionbreak

Marco war ihm über all die Jahre zu einem verlässlichen Freund geworden
und geblieben. Seinen beachtlichen, beruflichen Erfolg verdankte Louis
einzig ihm.

Beim Schritt von der Rolltreppe wäre Louis beinahe gestolpert.
Erschrocken hielt er sich am Gitarrenkoffer fest.

Ja, Marco hatte ihn herausgefordert. In vielerlei Hinsicht. Er hatte ihm
nicht nur eine stabile, zeitweise harte Ausbildung und Weiterbildung
geboten. Auch mit seiner Unterstützung in schulischen Belangen war er
standhaft an seiner Seite geblieben. Aber auch er, Louis, hatte alles
gegeben.

Er nagte an seiner Unterlippe und blickte in die Ferne.

Er hatte schnell gelernt. Zum ersten Mal hatte er ein berufliches Ziel
verfolgt. Und als schliesslich seine kulinarischen Eigenkreationen einen
unerwartet hohen Absatz fanden, war es geschafft. Louis war Teil von. Er
gehörte dazu. Er hatte etwas zu sagen. Die zusätzlich eingestellten
Köche wurden ihm kurz nach abgeschlossener Lehre unterstellt. Seine
Kochkunst hatte sich über die Monate herumgesprochen. Und eines Tages
hatte Marco mit seinem Lokal, nicht zuletzt durch ihn als Koch, den
angesehenen CCM, «Cucina Creativa Milano»-Preis gewonnen. Marco hatte
damals sogar mit dem Gedanken gespielt, Louis’ «Ossobuco-Kreation»
patentieren zu lassen. Louis hatte den italienischen Klassiker mit
seiner verblüffend eigenen Zubereitungsart neu erfunden und die Gäste
scharenweise angezogen.

Wie viel Abenteuer, wieviel Schönes ihn mit Marco verband. Marco hatte
ihm zu Selbstachtung verholfen. Diesem wahren Gefühl nach etwas
eigenständig Geschaffenen.

Dann war noch Vaters Bewunderung für ihn. Als Louis noch ein Kind war.
Doch irgend etwas war anders. Wie verloren er damals auch immer wieder
dagestanden hatte, wenn ihn Vaters Verhalten irritierte. Wenn er ihn aus
dem Nichts kritisierte und zeitnah Louis’ Befindlichkeit blitzartig
umdrehte, indem er ihn wieder aufbaute. Ihn über den Klee lobte. Ihn
stolz \emph{«mein kluger Sohn»} nannte. Ihn in die Luft warf, mit diesen
kräftigen Armen auffing, küsste und herzte.

In letzter Zeit kam ihm der Gedanke häufiger. Rätselhaft, überlegte er,
senkte den Kopf und kniff die Augen zusammen. Warum Vater nur so präsent
war? Vor allem, seit Louis gerade diese aufreibenden Kämpfe mit sich
ausfocht. Er hatte doch mit ihm abgeschlossen. Vor Jahren und endgültig.

\sectionbreak

\emph{Damals. Fribourg, sechs.}

\emph{Louis sitzt mit seinem Vater und Lucia im Auto. Maman ist zuhause
geblieben. Sie findet Vaters Idee unmöglich. Zu weit für ein paar
Stunden Spass. Zu teuer auch. Lucia zu klein. Nach drei Stunden
erreichen sie den Freizeitpark. Seine Schwester erwacht aus dem Schlaf.
Sie weint.}

\emph{«Ach, meine Prinzessin, lach mal, wir sind da,» klingt es laut und
gepresst fröhlich aus Vater. Er steigt hüpfend aus und hebt Lucia
überschwänglich aus dem Sitz. Sie schluchzt.}

\emph{«Die Windel, Papa, ist nass,» sagt sie und wischt sich über die Augen.}

\emph{Louis ist ausgestiegen und steht neben den beiden. Er fühlt sich ein
bisschen unnütz, als wäre er übrig geblieben und schaut unsicher zu.}

\emph{«Ja, ja, später,» ruft der Vater ihr zu. Er klingt genervt. Er hebt den
Rucksack aus dem Kofferraum. Louis überlegt, was er tun könnte.}

\emph{Lucia tut ihm leid. Er versteht sie ja. Sie ist stolz darauf, keine
Windeln mehr zu brauchen. Papa will, dass sie auf ihren Autofahrten
welche trägt. Das ist einfacher, hat er gesagt. Louis geht zu ihr hin.
Kaum daneben, umschlingt sie ihn mit ihren warmen, schmalen Ärmchen. Sie
beruhigt sich.}

\emph{«Sehr gut, Louis, das machst du sehr gut, mein Kumpel. Grosser Bruder,
nicht, Lucia?»}

\emph{Und der Vater starrt auf seine Videokamera, wechselt seine Haltung, geht
in die Knie.}

\emph{«Und jetzt, Kinder, lachen. Streckt die Arme in die Höhe, ruft hurra,
unser erstes Filmchen heute. Wir sind im Freizeitpark angekommen,
hurra.»}

\emph{Der Vater schreit die Worte hinaus und lacht dazu sehr laut. Lucia und
Louis ziehen die Lippen auseinander. Bei Lucia sieht es aus, als sei es
der Anfang eines erneuten Schluchzens. Louis hofft, dass sie nicht mehr
weint. Er streckt seine Arme in die Höhe. Dann senkt er sie nachdenklich
wieder, nimmt Lucias Hände und hebt sie hoch.}

\emph{Der Vater hüpft im Kreis. Jetzt steht er still, startet das Video und
kommentiert seine Aufnahme.}

\emph{«Da sind wir angekommen. Im tollsten Freizeitpark der Welt. Wie gefällt
es dir hier, Louis? Wohin geht es als erstes, mein Kumpel?» Seine
fröhliche Stimme klingt wie die des Moderators von Louis’
Lieblingssendung.}

\emph{Louis weiss nicht recht, was er sagen soll. Was dem Vater gefallen wird.
Lucia hält sich wieder an ihm fest. Sie wimmert leise und sagt erneut
etwas von dicker Windel.}

\emph{«Es gefällt mir gut, Papa. Vielleicht auf die lange Wasserrutsche?»}

\emph{Louis gibt sich viel Mühe, fröhlich zu klingen.}

\emph{Der Vater senkt die Kamera und fasst sich an den Kopf. Er verzieht das
Gesicht.}

\emph{«So geht das nicht, Kinder,» sagt er.}

\emph{Er klingt ein bisschen wütend. Louis zuckt zusammen.}

\emph{Dann kommt Vater auf Louis zu und hebt ihn hoch. Jetzt sieht er wieder
zufrieden aus, denkt Louis und ist beruhigt.}

\emph{«Was ist denn? Wir wollen doch fröhlich sein. Hab dich lieb wie nichts
anderes. Wir schaffen das zusammen. Männer können das. Du und ich, wir
sind die Coolsten, ja? Frauen sind manchmal ein bisschen kompliziert,»
dazu zeigt er auf Lucia und zwinkert Louis zu.}

\emph{Der Vater stellt ihn neben Lucia wieder ab. Sie steht breitbeinig da und
macht grosse Augen. Vater öffnet den Rucksack und sucht nach
Feuchttüchern. Louis fühlt sich besser. Er will gerne ein Mann sein wie
sein Vater. Der sieht nun auch wieder ziemlich lieb aus.}

\emph{Louis und Lucia bekommen vor den weiteren Video-Aufnahmen Gummibärchen.
Es läuft besser. Der Vater ist zufrieden.}

\sectionbreak

Louis setzte seine Flasche an. Das Wasser rann ihm kühl den Hals
hinunter. Damit fand er sich zurück auf dem Platz neben der Rolltreppe
vor dem Berner Bahnhof.

Die Sonne begann bereits wieder einzuheizen. Er hoffte, dass Marco den
Anruf entgegennahm. Sein Handy trug der stets mit sich. Ob er einen
Anruf mit unterdrückter Nummer annahm hingegen, war fraglich. Was würde
Louis erwarten?

Auf dem Platz war es hektisch und laut. Er schaute sich nach einer
ruhigen Ecke um. Etwas abseits fand er eine Bank und ging hinter einem
Abfalleimer in die Hocke.

Sein Herz begann beängstigend zu rasen. Er fühlte seinen Puls in den
Schläfen. Als er den Hörerknopf drücken wollte, rutschte ihm das Handy
um ein Haar aus der schweissnassen Hand.

Reiss dich zusammen, tu’s als Ciro. Er flüsterte die Worte vor sich hin.
Und wählte mit schwungvoller Überzeugung die Nummer.

Es klang, als würde jemand direkt neben seinem Ohr gurgeln. Dem folgte
ein Scheppern, ein Durcheinander an Stimmen und mittendrin, eingeklemmt,
ertönte ein gestresstes \emph{«Pronto?».} Dann klickte es und die
Leitung war tot. Piep, piep, piep. Louis senkte das Handy und ärgerte
sich.
